\documentclass[10pt,a4paper]{amsart}
\usepackage[margin=19mm]{geometry}
\usepackage{amsmath,amssymb,mathtools}
\usepackage[scr=rsfs,cal=euler]{mathalfa}
\usepackage{xfrac}
\usepackage[unicode,hidelinks,bookmarks=false]{hyperref}
\newcommand{\ie}{i.e.\ }
\newcommand{\cf}{cf.\ }
\newcommand{\resp}{respectively\ }
\newcommand{\Subsection}{\subsection}
\providecommand{\nobreakdash}{\nobreak\hbox{-}}
\setlength{\emergencystretch}{2em}
\multlinegap0pt
\usepackage{bm}
\usepackage{bbm}
\usepackage{dsfont}
\usepackage{xspace}
\usepackage{cancel}
\usepackage{mathtools}
\usepackage{amsthm}
\makeatletter
\@ifundefined{thm}{\newtheorem{thm}{Thm}}{}
\@ifundefined{lem}{\newtheorem{lem}[thm]{Lemma}}{}
\makeatother
\def\E{\mathbb E}
\def\P{\mathbb P}
\def\R{\mathbb R}
\renewcommand{\bar}{\overline}
\renewcommand{\d}{\hbox{d}}
\title[Rivers under Noise]{Rivers under Noise}
\author{Michael Scheutzow, Michael Grinfeld}
\date{Source: arXiv:2410.22207v3 (CC BY 4.0)}
\begin{document}
\maketitle

\noindent\emph{Source and licence.} Rivers under Noise. Version arXiv:2410.22207v3 (CC BY 4.0), distributed under the Creative Commons Attribution 4.0 International licence, as recorded in the arXiv licence metadata for this exact version and shown on the abstract page https://arxiv.org/abs/2410.22207v3. Full rights notice: licenses/arxiv-2410.22207-CC-BY-4.0.txt.\par\smallskip

\noindent\textit{Editorial scope.} Counted unit: the Lem whose source label is \texttt{lem2} in the article (source0.tex lines 949--958, source label \texttt{lem2}), delivered together with the article's own complete original proof of it (source0.tex lines 960--1065, one \texttt{proof} environment of 106 lines). No mathematics is written, repaired or re-derived: statement and proof are the article's own text line for line. Required context retained uncounted: kappe (lines 836--838); why (lines 869--871); 1ddiff (lines 1516--1551); solu (lines 1518--1520); ratio (lines 1530--1532). Every definition, display and label of those lines is retained unchanged. Omitted from this package and not redistributed: 1--835, 839--868, 872--948, 959--959, 1066--1515, 1521--1529, 1533--1617. Nothing inside a retained excerpt is omitted or altered: every retained body is delivered exactly as the article's lines are written, so no omission receipt of any kind accompanies this package. Citations: the retained excerpts carry no citation command at all, so no bibliography entry of the article is transcribed and no reference is redistributed. Reference adaptation: none was needed and none was made; every reference inside the delivered unit resolves inside it, so no unresolved reference or citation is produced. Printed slips kept literal: no printed word, symbol, hypothesis or number of the counted statement or of its proof was altered. Layout and class: the article's own class is \texttt{sn-jnl} with packages including color, graphicx, multirow, amsmath, amssymb, amsfonts, amsthm, mathrsfs, appendix, xcolor, textcomp, manyfoot, booktabs, algorithm; the wrapper is the lane's pinned \texttt{amsart} editorial wrapper at 10pt on A4, which restates the theorem-like environments the retained text needs and adds the article's own macro definitions that the retained text needs (\texttt{\textbackslash E}, \texttt{\textbackslash P}, \texttt{\textbackslash R}, \texttt{\textbackslash bar}, \texttt{\textbackslash d}), with extra wrapper packages (bm, bbm, dsfont, xspace, cancel, mathtools). The delivered numbering is the unit's own. Assets: no retained span contains a figure environment, a graphics-inclusion command, a PDF-page inclusion, a table or a TikZ picture; no image file is copied into this package, the package declares no asset dependency and no image file is redistributed with it. Licence: version arXiv:2410.22207v3 of this article is distributed under the Creative Commons Attribution 4.0 International licence, as recorded in the arXiv licence metadata for this exact version and shown on the abstract page https://arxiv.org/abs/2410.22207v3. A full rights notice is delivered at licenses/arxiv-2410.22207-CC-BY-4.0.txt. Characters: every retained excerpt is pure ASCII, so the delivered engine, XeLaTeX with its default Latin Modern text font, reports no missing character.
      \begin{equation}\label{kappe}
\kappa(s,x)\leq\int_\R \kappa(t,y)\,\P\big(X_{s,t}(x)\in\d y\big)
\end{equation}

\begin{equation}\label{why}
\d Y(t)=\big(Y(t) (Y(t) +t)-1     \big)\,\d t +\sigma\,\d W(t), \,t\geq s,\,Y(s)=x.
\end{equation}

\begin{lem}\label{1ddiff}
  Consider the scalar SDE
  \begin{equation}\label{solu}
\d Z(t) = b(Z(t))\, \d t + \sigma \,\d W(t), 
\end{equation}
where $\sigma>0$, $b:\R \to \R$ is locally Lipschitz continuous and $W$ is one-dimensional Brownian motion.
 
Fix some $c \in \R$ and define the {\em scale function}
$$
p(x):=\int_c^x\exp\Big\{ -\frac 2{\sigma^2}\int_c^\xi b(\zeta)\,\d \zeta\Big\}\,\d \xi,\;x \in \R.
$$
Let $-\infty\leq l < c<r \leq \infty$, let $Z$ be  the solution of \eqref{solu} with initial condition $x \in (l,r)$ and define  the exit time from the interval $(l,r)$ by $S:=\inf\{t \geq 0: Z(t) \notin (l,r)\}$. We define $p(r):=\lim_{y \uparrow r}p(y)$ if $r=\infty$ and similarly for $p(l)$. 
\begin{itemize}
\item [a)] If $p(l)>-\infty$ and $p(r)<\infty$ then $Z$ satisfies 
\begin{equation}\label{ratio}
\P\Big(\lim_{t \to S}Z(t)=l    \Big)=1-\P\Big(\lim_{t \to S}Z(t)=r    \Big)=\frac{p(r)-p(x)}{p(r)-p(l)}.
\end{equation}
(Note that the scale function depends on $c$ but the right hand side of \eqref{ratio} doesn't.)

\item[b)] If $p(l)>-\infty$ and $p(r)<\infty$, then
$$
\E\big[S\big]= -\int_l^x\big(p(x)-p(y)\big)\,m(\d y)+  \frac{p(x)-p(l)}{p(r)-p(l)}\int_l^r\big(p(r)-p(y)\big)\,m(\d y),
$$
where
$$
m(\d y):=\frac 2{p'(y)\sigma^2}\d y,\;y\in \R
$$
is the {\em speed measure} of $Z$.
\item[c)]
  Let
  $$
v(x):=\int_c^xp'(y)\int_c^y\frac 2{p'(z)\sigma^2}\,\d z\,\d y=\int_c^x\big(p(x)-p(y)\big)\,m(\d y).
$$
Then $\P (S=\infty)=1$ iff $v(l)=v(r)=\infty$. If $p(l)=-\infty$ and $v(r)<\infty$, then $S<\infty$ and $\lim_{t \uparrow S}Z(t)=\infty$ almost surely.  
  \end{itemize}
\end{lem}

\begin{lem}\label{lem2}
  For $x \in [-1,1]$, $Y(t),\,t \geq s$ with initial condition $Y(s)=x$ will exit the interval $[-1,1]$ in finite time, almost surely.   Let $p^+(s,x)$ be the probability that $Y(t)$, $t \geq s$ exits the interval $[-1,1]$ via 1. Then, for $z \in \R$,
  \begin{equation}\label{limm}
\lim_{s \to \infty}p^+\Big(s,z\frac {\sigma}{\sqrt{2s}}\Big)= \Phi(z):=\frac 1{\sqrt{2\pi}}\int_{-\infty}^z\exp\Big\{-\frac{y^2}2 \Big\}\,\d y.
\end{equation}
Further, for $\alpha \in [0,1/2)$,
$$
p^+\big(s,-s^{-\alpha}\big) \leq \exp\Big\{-\frac 1{\sigma^2}s^{1-2\alpha}\big(1+o(1)\big)    \Big\},\;s \to \infty.
$$
\end{lem}

\begin{proof}
  Instead of $Y$, we first consider the solution $Z_A$ of
 \begin{equation}\label{ZZZ}
\d Z_A(t)=\big(sZ_A(t)+A\big)\,\d t +\sigma \d W(t),\,t \geq s;\, Z_A(s)=x \in [-1,1],
\end{equation}
where $A \in \R$ may depend on $s$ (but neither on $t$ nor on $Z$).
Let
$$
q_A(x):=\int_0^x\exp \Big\{-\frac 2{\sigma^2} \int_0^y\big(us+A\big)\,\d u   \Big\}\,\d y
$$
be the scale function of $Z_A$. Then
\begin{align*}
q_A(x)&=\int_0^x \exp \Big\{-\frac 2{\sigma^2} \Big(\frac{sy^2} 2+Ay\Big)   \Big\}\,\d y\\
&=\int_{\frac{2A}{\sigma \sqrt{2s}}}^{\frac {\sqrt{2s}x}\sigma + \frac{2A}{\sigma\sqrt{2s}}} \exp\Big\{-\frac 12 v^2\Big\}\,\d v \cdot \frac \sigma{\sqrt{2s}}  \exp\Big\{\frac {A^2}{\sigma^2 s}\Big\}.
\end{align*}
By Lemma \ref{1ddiff}a), the probability $p_A(s,x)$ that $Z_A$ exits the interval $[-1,1]$ via 1 equals
\begin{equation}\label{frac}
\frac{q_A(x)-q_A(-1)}{q_A(1)-q_A(-1)}=\frac{\Phi\Big( \frac {\sqrt{2s}x}\sigma +\frac{2A}{\sigma \sqrt{2s}}\Big) - \Phi\Big( -\frac {\sqrt{2s}}\sigma +\frac{2A}{\sigma \sqrt{2s}}\Big)}
{\Phi\Big( \frac {\sqrt{2s}}\sigma +\frac{2A}{\sigma \sqrt{2s}}\Big) - \Phi\Big( -\frac {\sqrt{2s}}\sigma +\frac{2A}{\sigma \sqrt{2s}}\Big)}.
\end{equation}
For given $z \in \R$, define $x=z\frac \sigma{\sqrt{2s}}$. Then, by \eqref{frac},
\begin{equation}\label{assame}
\lim_{s \to \infty} p_A\Big(s, z\frac \sigma{\sqrt{2s}}\Big)=\Phi(z),
\end{equation}
provided that $A=A(s)=o(\sqrt{s})$.

Next, we compare  $Y$ and, for specific functions $A$, $Z_A$  both with the same initial condition $x\in [-1,1]$ at time $s$.
For $A=A(s)=o\big(s^{1/2}\big)$ let
$$
T(s,x,A):=\inf\{t \geq s: |Z_A(t,x)|=1\}
$$ denote the first time after $s$ when  $Z_A$ starting at $Z_A(s)=x\in[-1,1]$ hits $\{-1,1\}$. We will show that, for every such function $A=o\big(s^{1/2}\big)$,
\begin{equation}\label{zuzeigen}
  \zeta_A:=\limsup_{s \to \infty}\Big(\sup_{x \in [-1,1]}\E \big[T(s,x,A)\big]-s\Big)<\infty.
\end{equation}

Once we have established \eqref{zuzeigen}, we apply it to 
$\bar A(s):=s^\alpha$, $\underline A(s):=-s^\alpha-1 $ for some fixed value of $\alpha \in (0,1/2)$.  We estimate the drift function $b(t,y)$ of $Y$ from above
by $b(t,y)=y^2-1+yt= y^2-1+ys+y(t-s) \leq ys + \bar A(s)$ %whenever $t \in [s,s+s^\alpha]$, and $y \in [-1,1]$
and, from below, by $b(t,y)=y^2-1+yt= y^2-1+ys+y(t-s) \geq ys + \underline A(s)$ whenever $t \in [s,s+s^\alpha]$, and $y \in [-1,1]$.

Fix $s >0$ and $z \in \R$ such that $x:=z\frac{\sigma}{\sqrt{2s}}\in[-1,1]$ and consider  $Z_{\bar A}$ and $Z_{\underline A}$ with initial condition $x$ at time $s$.
%be the solution of \eqref{ZZZ} with initial condition $x$ %$z\frac \sigma{\sqrt{2s}}$
%and $A$ replaced by $\bar A$ and $\underline A$, respectively.
By comparison, we have $Z_{\underline A}(t)\leq Y(t) \leq Z_{\overline A}(t)$ for $t \in [s,s+s^\alpha]$ up to the minimum of the exit times $T(s,x,\bar A)$ and  $T(s,x,\underline A)$ of
$Z_{\overline A}$ and $Z_{\underline A}$ from $[-1,1]$, where $Y$ denotes the solution of \eqref{why} with the same initial condition $x$.  We saw that, as $s \to \infty$, the probability of exiting the
interval $[-1,1]$ via 1 is asymptotically the same for   $Z_{\underline A}$ and for $Z_{\bar A}$. In order to ensure that this is also true for $Y$ we have to show that the exit times from $[-1,1]$
of these processes are at most $s+s^{\alpha}$ with probabilities converging to 1 as $s \to \infty$. We have
\begin{align*}
  p^+(s,x)&\geq \P\big(Z_{\underline{A}}\mbox{ exits }[-1,1] \mbox{ via } 1 \mbox{ before or at time } s+s^{\alpha}\big)\\
  &=\P\big( Z_{\underline{A}}\mbox{ exits }[-1,1] \mbox{ via } 1\big)-  \P\big(Z_{\underline{A}}\mbox{ exits }[-1,1] \mbox{ via } 1 \mbox{ after time }s+s^{\alpha}\big).
\end{align*}
Define $p^-(s,x)$ as the probability that $Y$ starting at $x\in [-1,1]$ at time $s$ exits $[-1,1]$ via -1 (clearly, $p^+(s,x)+p^-(s,x)\leq 1$ but it is not yet clear that there is equality). Then, analogously,
\begin{align*}
  p^-(s,x)&\geq \P\big(Z_{\bar{A}}\mbox{ exits }[-1,1] \mbox{ via } -1 \mbox{ before or at time } s+s^{\alpha}\big)\\
  &=\P\big( Z_{\bar{A}}\mbox{ exits }[-1,1] \mbox{ via } -1\big)-  \P\big(Z_{\bar{A}}\mbox{ exits }[-1,1] \mbox{ via } -1 \mbox{ after time }s+s^{\alpha}\big).
\end{align*}  
By Markov's inequality, \eqref{assame} and \eqref{zuzeigen}, we have
$$
 p^+(s,x)+ p^-(s,x)\geq 1+o(1)-\frac{\zeta_{\bar A}}{s^{\alpha}}- \frac{\zeta_{\underline A}}{s^{\alpha}}\to 1,\;s \to \infty,
 $$
 and the same argument as in \eqref{kappe} shows that we actually have $ p^+(s,x)+ p^-(s,x)=1$ for every $|x|\leq 1$ and $s$, so the very first statement in the lemma follows. Further,
 \begin{align*}
   1-\P\big( Z_{\underline{A}}\mbox{ exits }[-1,1] \mbox{ via } -1\big)+o(1)&\geq 1-p^-\big(s,z\frac {\sigma}{\sqrt{2s}}\big)\\
   &=p^+\big(s,z\frac {\sigma}{\sqrt{2s}}\big)\geq \P\big( Z_{\bar{A}}\mbox{ exits }[-1,1] \mbox{ via } 1\big)+o(1),
\end{align*}
and therefore, by \eqref{assame}, \eqref{limm} follows.



%$$
%p^+(s,x)\leq \P\big( Z_{\underline A} \mbox{ exits via 1 before } s+s^{-\alpha}\big)\leq \P\big( Z_{\underline A} \mbox{ exits via 1}\big) - \P\big( Z_{\underline A} 
%$$

It remains to prove \eqref{zuzeigen}. Define
$$
V(t):=\Big(Z_A(t,x)+\frac As\Big)^2,\,t \geq s.
$$
and assume that $s$ is so large that $\frac As \in (-1,1)$. Then, by It\^o's formula, for $t \geq s$,
$$
V(t)=\Big(x+\frac As\Big)^2+2s\int_s^t V(u)\,\d u+M_t+\sigma^2 (t-s),
$$
where $M_t$, $t \geq s$ is a continuous local martingale satisfying $M_s=0$. Let $\tau$ be the first time after $s$ when $V$ attains the value 4. Then, since $M_{t\wedge \tau}$ is a martingale starting at 0 and since $V \geq 0$, we obtain
$$
4 \geq \Big(x+\frac As\Big)^2 + \sigma^2 \Big(\E [\tau]-s\Big),
$$
so
$$
\E [\tau] \leq s+ \frac 4{\sigma^2}.  
$$
This implies $\E \big[T(s,x,A)\big]-s \leq \E [\tau]-s \leq \frac 4{\sigma^2}$ for all $x \in [-1,1]$. %provided that $s$ is so large that $\frac As \in [-1,1]$.
This completes the proof of the first part of the lemma.

The final claim in the lemma   follows easily from \eqref{frac} :
\begin{align*}
p^+\big(s,-s^{-\alpha}\big)&\leq p_0(s,-s^{-\alpha}\big)=\frac{\Phi\Big(-\frac1\sigma \sqrt{2s}s^{-\alpha}\Big)-\Phi\Big(-\frac1\sigma \sqrt{2s}\Big)}{\Phi\Big(\frac1\sigma \sqrt{2s}\Big)-\Phi\Big(-\frac1\sigma \sqrt{2s}\Big)}\\
&\sim  \Phi\Big(-\frac1\sigma \sqrt{2}s^{\frac 12-\alpha}\Big)  \leq\exp\Big\{-\frac 1{\sigma^2}s^{1-2\alpha}\big(1+o(1)\big)    \Big\}.
\end{align*}


%We regard $s$ and $A$ as fixed and denote the expected value of $T(s,x,A)$ by $M(x)$.
%By \cite[(5.55)]{KS}, we have
%$$
%M(x)=\frac {q(1)-q(x)}{q(1)-q(-1)}\int_{-1}^x \big(q(y)-q(-1)\big)\,m(\d y)     + \frac {q(x)-q(-1)}{q(1)-q(-1)}\int_x^1\big(q(1)-q(y)\big)\, m(\d y), 
%$$
%where $m(\d y)=\frac 2{q'(y)\sigma^2}\,\d y$ is the {\em speed measure} of $Z_A$. Then ...
\end{proof}

\end{document}
